% =====================================================================
%  Section: Introduction
% =====================================================================
\section{Introduction}
\label{sec:intro}

Classical paintings---whether oil on canvas, ink on silk, or pigment on
fresco---are cultural artifacts that encode centuries of human visual
imagination~\cite{lowenthal2015heritage}. From van Gogh's swirling
\textit{Starry Night} to Zhang Zeduan's panoramic \textit{Along the River During
the Qingming Festival}, these works imply motion, narrative, and spatial depth
that their static medium can only suggest. Animating such paintings as
explorable videos has therefore emerged as a compelling research goal at the
intersection of computer vision, generative modeling, and digital cultural
heritage~\cite{zhang2024artgallery,dong2024enhancingpainting}.

Recent years have witnessed remarkable progress on two fronts. On one hand,
image-to-video (I2V) diffusion models---including Stable Video
Diffusion~\cite{blattmann2023svd}, Wan2.2~\cite{wan2025},
LTX-Video~\cite{ltxvideo2024}, and CogVideoX~\cite{yang2025cogvideox}---can
animate a single photograph into a short clip with plausible dynamics. Combined
with ControlNet~\cite{zhang2023controlnet} and its video
extensions~\cite{wang2024easycontrol,peng2024controlnext}, these models achieve
impressive per-frame visual quality and controllable motion. Yet when applied
directly to oil paintings or ink-wash artworks, they suffer from two fundamental
shortcomings. First, they operate entirely in 2D pixel space: the model has no
geometric understanding of object depth ordering, self-occlusion, or relative
spatial layout, and consequently cannot synthesize novel viewpoints or maintain
consistent 3D structure across frames. Second, long-range temporal coherence
degrades rapidly beyond 2--4 seconds, as the latent-space interpolation lacks
explicit geometric constraints to anchor the motion trajectory.

On the other hand, 3D reconstruction and rendering methods---Neural Radiance
Fields (NeRF)~\cite{mildenhall2020nerf}, 3D Gaussian
Splatting~(3DGS)~\cite{kerbl20233dgs}, and 2D Gaussian
Splatting~(2DGS)~\cite{huang20242dgs}---can reconstruct explorable 3D scenes
from sparse views. Monocular feed-forward variants such as
TripoSplat~\cite{wang2024triposplat} and DreamGaussian~\cite{tang2024dreamgaussian}
can even lift a single image to a 3D representation. However, these
reconstruction pipelines are optimized for photorealistic fidelity;
when applied to paintings, the rendering process replaces expressive
brushstrokes with geometric artifacts, flattens the impasto texture, and bleeds
semantic color boundaries under novel-view projection.

The core tension is this: \textit{3D scene structure and artistic style
fidelity} are governed by incompatible optimization objectives, and no single
model family---neither pure diffusion nor pure 3D reconstruction---satisfactorily
achieves both. This tension becomes especially acute when the animation task
requires \textbf{object-level manipulation}: moving a painted figure from one
compositional position to another while preserving its painterly identity and
filling the resulting gaps with plausible background content.

\smallskip
\noindent\textbf{Our framework: Canvas3D.}
We propose \textbf{Canvas3D}, a framework that resolves this tension by
separating the problem into four interacting yet independently optimized
sub-systems, each specialized for one aspect of the painting-to-video
pipeline (Figure~\ref{fig:overview}). \textbf{PaintParse} decomposes the input
painting into foreground objects and a background layer using a fine-tuned
YOLO-family instance segmentation model, augmented with monocular depth
estimation from Depth Anything V2 for metric-scale layer ordering.
\textbf{MonoSplat} lifts each segmented foreground region into a high-fidelity
3DGS asset via feed-forward TripoSplat reconstruction, assembling an
object-level model library with depth-derived metric rescaling.
\textbf{SketchComp} accepts hand-drawn line-art sketches specifying target
object positions and spatial layout, solves a contour-based rigid registration
to place each 3DGS object onto its sketch, and fills disocclusion gaps through a
depth-aware inpainting module that preserves background artistic texture.
\textbf{ArtiFlow} renders the composited 3D scene into multi-view keyframes and
feeds them into a ComfyUI-based video synthesis pipeline combining Wan2.2 or
LTX2.3 with dual ControlNet modules conditioned on depth and HED edge maps,
producing temporally smooth animations with pixel-precise object motion governed
by explicit 3D geometry rather than learned heuristics.

This decomposition enables a clean separation of concerns: PaintParse handles
perceptual grouping, MonoSplat handles geometric structure, SketchComp handles
artistic intent, and ArtiFlow handles temporal rendering and style refinement.
The result is a system that supports \textit{arbitrary object repositioning},
\textit{multi-view explorability}, and \textit{style-anchored animation}---
capabilities that no prior single-modality pipeline offers.

\smallskip
\noindent\textbf{Contributions.} Our contributions are:
\begin{enumerate}
    \item \textbf{The Canvas3D framework}, integrating PaintParse (art-domain
    segmentation), MonoSplat (per-object 3DGS lifting), SketchComp (line-art
    composition with depth inpainting), and ArtiFlow (dual-ControlNet video
    synthesis) into the first end-to-end pipeline for 3D-aware, artistically
    controllable painting animation.
    \item \textbf{An object-level 3DGS asset pipeline} (MonoSplat) that combines
    YOLO-based instance segmentation with TripoSplat monocular lifting,
    producing a reusable library of 3D foreground objects from a single
    painting.
    \item \textbf{A sketch-conditioned 3D scene composition algorithm}
    (SketchComp) with contour-based rigid registration and a depth-aware color
    inpainting module that fills disocclusion gaps while preserving the
    background's artistic texture.
    \item \textbf{A geometry-grounded video synthesis workflow} (ArtiFlow)
    combining Wan2.2/LTX2.3 with dual ControlNet (depth + HED edge) for
    explicit 3D-structure-anchored animation with per-object motion control.
    \item \textbf{Comprehensive evaluation} on a curated benchmark of 50
    classical oil paintings and ink-wash works, achieving state-of-the-art
    style fidelity (CLIP-I: 0.86), perceptual quality (LPIPS: 0.28), and
    temporal coherence (FID-VID: 18.3), with user studies confirming 82\%
    preference for artistic authenticity and 88\% for compositional intent
    alignment.
\end{enumerate}
